% !TeX root = ../main.tex
\begin{abstract}
Encrypted transport protocols such as QUIC hide application payloads, yet flow-level
metadata---packet sizes, directions, and inter-packet times---remain observable to
network operators and passive adversaries. Deep learning classifiers can exploit this
metadata to infer which application produced a flow. We study \emph{deterministic}
metadata obfuscation as a countermeasure on real CESNET-QUIC22 backbone traces.
Using a temporally held-out test week (49{,}305 flows, 64 application classes), we
train and freeze two attack architectures---a masked Transformer encoder and a
CNN-BiLSTM---and evaluate seven obfuscation policies spanning timing jitter, linear
size padding, and MTU padding, with manifest-derived bandwidth and latency overhead.
On a masked Transformer attacker, clean traffic yields 74.4\% macro F1 (77.8\%
accuracy). Moderate Laplace jitter with only 11\,ms mean latency overhead
(\emph{jitter\_low}) reduces macro F1 by 1.7 percentage points while preserving
76.8\% accuracy; the degradation is statistically significant (McNemar
$p \approx 5.8\times10^{-22}$) but small in magnitude. Strong jitter and MTU padding
achieve near-random performance at very different cost points. Transformers remain
more robust than CNN-BiLSTMs under jitter, while both architectures collapse under
MTU obfuscation. We release a reproducible pipeline and formal Pareto analysis to
guide proxy deployment.
\end{abstract}

\begin{IEEEkeywords}
QUIC, traffic classification, metadata obfuscation, privacy, Pareto frontier, deep learning
\end{IEEEkeywords}

%==============================================================================
